% LuaLaTeX でコンパイルしてください(Overleaf: Menu → Compiler → LuaLaTeX)
\documentclass[a4paper,11pt]{ltjsarticle}

\usepackage[margin=25mm]{geometry}
\usepackage{amsmath,amssymb,bm}
\usepackage{booktabs}
\usepackage[most]{tcolorbox}
\usepackage[hidelinks]{hyperref}

\newcommand{\Win}{W_{\mathrm{in}}}
\newcommand{\Wout}{W_{\mathrm{out}}}
\newcommand{\yt}{\bm{y}_{\mathrm{teach}}}
\newcommand{\T}{^{\top}}
\newcommand{\tr}{\operatorname{tr}}

\newtcolorbox{keybox}[1][]{colback=blue!4,colframe=blue!50!black,
  fonttitle=\bfseries,arc=2pt,boxrule=0.6pt,#1}

\title{リザバーコンピューティングの学習と\\最適制御的なリードアウト設計}
\date{}

\begin{document}
\maketitle

%===========================================================
\section{リザバーコンピューティングの学習}
%===========================================================

リザバーコンピューティング(RC)は,ネットワーク本体は学習せず,
出力層(リードアウト)だけを学習する手法である.

\subsection{構造}

\begin{align}
  \bm{x}(t+1) &= \tanh\!\bigl(W\bm{x}(t) + \Win\bm{u}(t)\bigr)
  && \text{(リザバー: $W,\Win$ はランダムに固定)} \\
  \bm{y}(t) &= \Wout\,\bm{x}(t)
  && \text{(リードアウト: ここだけ学習)}
\end{align}

\begin{center}
\begin{tabular}{lll}
\toprule
対象 & 学習するか & 備考 \\
\midrule
入力重み $\Win$ & しない & ランダム固定 \\
内部結合 $W$ & しない & スペクトル半径のみ調整 \\
出力重み $\Wout$ & \textbf{する} & リッジ回帰など \\
ハイパーパラメータ & しない & 検証データで選択 \\
\bottomrule
\end{tabular}
\end{center}

\subsection{リッジ回帰による学習}

訓練入力をリザバーに流し,各時刻の状態を列に並べた行列
$X = [\bm{x}(1),\dots,\bm{x}(T)]$ と,教師信号の行列
$Y = [\yt(1),\dots,\yt(T)]$ を作る
(最初の数百ステップは初期値の影響が残るため,ウォッシュアウトとして捨てる).

\begin{keybox}[title=リッジ回帰の解]
\begin{equation}
  \Wout = Y X\T \bigl(X X\T + \lambda I\bigr)^{-1}
  \label{eq:ridge}
\end{equation}
\end{keybox}

勾配降下も BPTT も不要で,RNN 学習における勾配消失・爆発の問題を回避できる.

\subsection{ランダムなリザバーで動く理由}

\begin{itemize}
  \item \textbf{エコーステート性}:状態が初期値ではなく入力履歴だけで決まる性質.
        実用上は $W$ のスペクトル半径 $\rho(W)$ を $1$ 弱($0.9$--$0.99$ 程度)にスケーリングする.
  \item \textbf{高次元への非線形展開}:入力履歴が高次元状態空間に非線形に埋め込まれ,
        線形読み出しだけで複雑な関数を近似できる(カーネル法に近い発想).
  \item \textbf{メモリと非線形性のトレードオフ}:$\rho(W)$,入力スケール,リーク率などで調整する.
\end{itemize}

%===========================================================
\section{リッジ回帰は既に LQR 型の目的関数である}
%===========================================================

制御理論で極配置 $\bm{u} = -K\bm{x}$ の代わりに目的関数を最小化する $\bm{u}$ を考えたのと同様に,
リードアウト学習の目的関数を LQR 風に一般化する.

\begin{keybox}[title=一般化した目的関数]
\begin{equation}
  J(\Wout) = \sum_{t} \bm{e}(t)\T Q\, \bm{e}(t)
           + \tr\!\bigl(\Wout R\, \Wout\T\bigr),
  \qquad
  \bm{e}(t) = \yt(t) - \Wout\bm{x}(t)
  \label{eq:J}
\end{equation}
\end{keybox}

\begin{itemize}
  \item $Q \succ 0$:誤差の重み(LQR の状態コストに相当)
  \item $R \succ 0$:重みの大きさへのペナルティ(LQR の制御コストに相当)
\end{itemize}

通常のリッジ回帰は $Q = I$, $R = \lambda I$ の特殊ケースである:
\begin{equation}
  J = \underbrace{\sum_t \bigl\|\yt(t) - \Wout\bm{x}(t)\bigr\|^2}_{\text{状態コスト($Q$ 相当)}}
    + \underbrace{\lambda\,\|\Wout\|_F^2}_{\text{制御コスト($R$ 相当)}} .
\end{equation}

\subsection{最適解:Sylvester 方程式}

$R$ を対称として $\partial J / \partial \Wout = 0$ とおくと
\begin{equation}
  -2\,Q\bigl(Y - \Wout X\bigr)X\T + 2\,\Wout R = 0
\end{equation}
より,次の Sylvester 方程式を得る.

\begin{keybox}[title=最適リードアウトの条件]
\begin{equation}
  Q\,\Wout\,(X X\T) + \Wout R = Q\, Y X\T
  \label{eq:sylvester}
\end{equation}
\end{keybox}

これは \texttt{scipy.linalg.solve\_sylvester} で直接解ける.
$Q = I,\ R = \lambda I$ とすれば式~\eqref{eq:ridge} に戻る.

\subsection{対角 $Q$ の場合:出力ごとの実効正則化}

$Q = \mathrm{diag}(q_1,\dots,q_m)$,$R = \lambda I$ のとき,
$\Wout$ の第 $i$ 行 $\bm{w}_i$ と $Y$ の第 $i$ 行 $\bm{y}_i$ について問題は分解でき,
\begin{equation}
  \bm{w}_i = \bm{y}_i X\T \Bigl(X X\T + \frac{\lambda}{q_i} I\Bigr)^{-1}
\end{equation}
となる.すなわち実効的な正則化強度は $\lambda / q_i$ である.

\begin{itemize}
  \item $q_i$ 大:その出力は忠実にフィットする
  \item $q_i$ 小:その出力は多少ずれてもよいから重みを小さくする(汎化重視)
\end{itemize}

%===========================================================
\section{$R$ の設計:状態空間のどの方向を信用するか}
%===========================================================

$X$ の特異値分解を $X = U \Sigma V\T$($\Sigma = \mathrm{diag}(\sigma_1,\sigma_2,\dots)$)とする.
リッジ回帰の当てはめ値は
\begin{equation}
  \hat{Y} = \Wout X = Y \sum_i \bm{v}_i
    \underbrace{\frac{\sigma_i^2}{\sigma_i^2 + \lambda}}_{\text{縮小率}} \bm{v}_i\T
\end{equation}
と書け,各主成分方向の寄与を縮小率 $\sigma_i^2/(\sigma_i^2+\lambda)$ で抑えている.

\begin{itemize}
  \item 分散が大きい方向($\sigma_i^2 \gg \lambda$):ほぼそのまま使う(縮小率 $\approx 1$)
  \item 分散が小さい方向($\sigma_i^2 \ll \lambda$):強く抑える(縮小率 $\approx 0$)
\end{itemize}

分散の小さい方向はノイズに敏感で,
\textbf{訓練データには効くが予測では裏切る}方向である.そこで
\begin{itemize}
  \item $R = \Gamma\T\Gamma$ を「低分散方向ほど強く罰する」よう非等方に設計する(一般化 Tikhonov 正則化)
  \item 極端には低分散方向を完全に切り捨てる(打ち切り SVD/主成分回帰)
\end{itemize}
とすれば,「訓練精度を捨てて予測精度を取る」を方向ごとに制御できる.

\paragraph{補足(推定と制御の双対性)}
リッジ回帰は,$\Wout$ の各行に事前分布 $\mathcal{N}(\bm{0},\,R^{-1})$ を置いた MAP 推定と解釈できる.
これはカルマンフィルタと LQR の双対性とも整合する.

%===========================================================
\section{LQR の類推が本当に効くのは閉ループ予測}
%===========================================================

自律生成モード(自分の出力を入力にフィードバックする)では,状態の時間発展が変わる:
\begin{align}
  \text{開ループ(訓練時,teacher forcing)}:\quad
  & \bm{x}(t+1) = \tanh\!\bigl(W\bm{x}(t) + \Win\yt(t)\bigr) \\
  \text{閉ループ(予測時,自由走行)}:\quad
  & \bm{x}(t+1) = \tanh\!\bigl((W + \Win\Wout)\,\bm{x}(t)\bigr)
  \label{eq:closed}
\end{align}

式~\eqref{eq:closed} を線形化すると $\bm{x}(t+1) \approx D(t)\,(W + \Win \Wout)\,\bm{x}(t)$
($D(t)$ は $\tanh$ の微分を並べた対角行列)となり,線形状態フィードバック系
$\bm{x}(t+1) = (A - BK)\bm{x}(t)$ と次のように対応する.

\begin{center}
\begin{tabular}{ccc}
\toprule
制御理論 & & リザバー \\
\midrule
$A$ & $\leftrightarrow$ & $W$ \\
$B$ & $\leftrightarrow$ & $\Win$ \\
$-K$ & $\leftrightarrow$ & $\Wout$ \\
\bottomrule
\end{tabular}
\end{center}

\begin{keybox}[title=要点]
$\Wout$ は\textbf{状態フィードバックゲインそのもの}である.
訓練時の 1 ステップ誤差を最小化しても閉ループの安定性は保証されない.
「訓練精度と予測精度の乖離」の正体は,\textbf{開ループ最適 $\neq$ 閉ループ安定}という
制御でおなじみの問題である.
\end{keybox}

\subsection{既知の処方箋と $Q, R$ との対応}

\paragraph{(1) ノイズ注入学習}
訓練時の状態に平均 $\bm{0}$,共分散 $\Sigma_n$ のノイズ $N$ を加えて回帰すると,
\begin{equation}
  \mathbb{E}\bigl[(X+N)(X+N)\T\bigr] = X X\T + T\,\Sigma_n
\end{equation}
より,期待値の意味で $R = T\,\Sigma_n$ の一般化 Tikhonov 正則化と等価になる(Bishop, 1995).
\textbf{ノイズの共分散がそのまま $R$ の役割を果たし},
フィードバック誤差が蓄積しても発散しにくい $\Wout$ に寄る.

\paragraph{(2) FORCE 学習}
RLS(再帰最小二乗法)で誤差を常に小さく保ちながら,閉ループのまま $\Wout$ を更新する.
オンライン LQG に近い考え方である.

\paragraph{(3) 多段先予測コスト(MPC のホライズンコストの類推)}
\begin{keybox}[title=多段先予測コスト]
\begin{equation}
  J(\Wout) = \sum_{t}\sum_{k=1}^{H} \gamma^{k}\,
    \bigl\|\hat{\bm{y}}(t+k \mid t) - \yt(t+k)\bigr\|^2
\end{equation}
\end{keybox}
ここで $\hat{\bm{y}}(t+k\mid t)$ は時刻 $t$ から式~\eqref{eq:closed} で $k$ ステップ自由走行させた予測値,
$H$ は予測ホライズン,$\gamma \in (0,1]$ は割引率である.

閉ループのロールアウトを含むため閉形式では解けず,
ロールアウトを通した勾配降下(微分可能シミュレーション)で解く.
「1 ステップの訓練精度を明示的に犠牲にして $k$ ステップ先の予測を最適化する」という,
最も直接的な定式化である.

%===========================================================
\section{進め方の案}
%===========================================================

\begin{enumerate}
  \item 標準 ESN + リッジ回帰でベースラインを作る(例:Lorenz 系の自律生成).
        「訓練誤差は小さいのに自由走行で軌道が崩れる」現象を再現する.
  \item $Q, R$ 付き Sylvester 版(式~\eqref{eq:sylvester})とノイズ注入版で,
        閉ループ性能の変化を比較する($\lambda$ の掃引 vs.\ $R$ の形状設計).
  \item 多段先予測コストを勾配法で解き,閉形式解との差を見る.
  \item 発展:磁石シミュレータ(磁気双極子系の非線形ダイナミクス)を
        物理リザバーに見立てて同じ実験を行う.
\end{enumerate}

\end{document}
